% GENERATED FILE. Edit canonical lessons, then run npm run generate:books.
% canonical-source-hash: fnv1a64:177f71a1f8ad0205
% canonical-lessons: RU-C367-svoboda, RU-C367-ravenstvo, RU-C367-bednost, RU-C367-bogatstvo, RU-C367-bezrabotica

\chapter{Freedom and equality}
\label{ch:ru-freedom-and-equality}
\hlchaptermodality{367}

\begin{chapteropening}
\textbf{By the end of this chapter:} \emph{I can talk about freedom and equality, poverty and wealth, and unemployment.}

Complete the last lesson of chapter 367: I can talk about freedom and equality, poverty and wealth, and unemployment.
\end{chapteropening}

\section[\texorpdfstring{\ru{свобода} (svobóda)}{свобода (svobóda)}]{\ru{свобода} (svobóda) — freedom}

\label{lesson:RU-C367-svoboda}



\begin{quote}
Before the new one: say the Russian for a protest, then the Russian for to protest.
\end{quote}

\subsection*{You'll want to know: \ru{свобода}}
\textbf{\ru{свобода}} (\emph{svobóda}) — \textquotedblleft{}freedom\textquotedblright{}.

\begin{quote}
\ru{Для} \ru{меня} \ru{свобода} \ru{важнее} \ru{денег}.
\end{quote}

\begin{quote}
\emph{Dlya menyá svobóda vazhnéye déneg.}
\end{quote}

\begin{quote}
For me, freedom matters more than money.
\end{quote}

Feminine. \emph{\ru{Свобода} \ru{слова}} — freedom of speech. \emph{\ru{На} \ru{свободе}} — at large.

\begin{cousinweb}
\textbf{\ru{свобода}} is inherited through Old East Slavic \emph{svoboda} from Proto-Slavic \emph{*svoboda} \textquotedblleft{}freedom\textquotedblright{} (EtymDB), usually traced to the root of \emph{\ru{свой}} \textquotedblleft{}one’s own\textquotedblright{}.
\end{cousinweb}

\begin{grammarlens}[title={What you've built}]
One more word for telling, explaining and describing at B1.
\end{grammarlens}

\subsection*{Your turn}
\begin{itemize}
\raggedright
  \item \emph{Say it:} \emph{svobóda}
  \item \emph{Say it:} \emph{svobóda}, once more
  \item \emph{Say it:} a sentence of your own with \emph{svobóda}
  \item \emph{Recall:} say the Russian for a protest, then the Russian for to protest, then say \emph{svobóda} again
\end{itemize}

\begin{tcolorbox}[breakable,colback=teal!4,colframe=teal!35!black,title={Before you move on}]
What does \ru{свобода} mean? (\textquotedblleft{}Freedom\textquotedblright{}.) Say it once more.
\end{tcolorbox}
\section[\texorpdfstring{\ru{равенство} (rávenstvo)}{равенство (rávenstvo)}]{\ru{равенство} (rávenstvo) — equality}

\label{lesson:RU-C367-ravenstvo}



\begin{quote}
Before the new one: say the Russian for to protest, then the Russian for freedom.
\end{quote}

\subsection*{You'll want to know: \ru{равенство}}
\textbf{\ru{равенство}} (\emph{rávenstvo}) — \textquotedblleft{}equality\textquotedblright{}.

\begin{quote}
\ru{Мы} \ru{верим} \ru{в} \ru{равенство} \ru{всех} \ru{людей}.
\end{quote}

\begin{quote}
\emph{My vérim v rávenstvo vsekh lyudéy.}
\end{quote}

\begin{quote}
We believe in the equality of all people.
\end{quote}

Neuter. \emph{\ru{Равенство} \ru{мужчин} \ru{и} \ru{женщин}} — equality of men and women.

\begin{cousinweb}
\textbf{\ru{равенство}} is built on \emph{\ru{равный}} \textquotedblleft{}equal\textquotedblright{}, the Church Slavonic twin of native \emph{\ru{ровный}} \textquotedblleft{}even, level\textquotedblright{}; both go back to Proto-Slavic \emph{*orvĭnŭ}. The \emph{\ru{ра}-} is the Church Slavonic shape, \emph{\ru{ро}-} the native one.
\end{cousinweb}

\begin{grammarlens}[title={What you've built}]
One more word for telling, explaining and describing at B1.
\end{grammarlens}

\subsection*{Your turn}
\begin{itemize}
\raggedright
  \item \emph{Say it:} \emph{rávenstvo}
  \item \emph{Say it:} \emph{rávenstvo}, once more
  \item \emph{Say it:} a sentence of your own with \emph{rávenstvo}
  \item \emph{Recall:} say the Russian for to protest, then the Russian for freedom, then say \emph{rávenstvo} again
\end{itemize}

\begin{tcolorbox}[breakable,colback=teal!4,colframe=teal!35!black,title={Before you move on}]
What does \ru{равенство} mean? (\textquotedblleft{}Equality\textquotedblright{}.) Say it once more.
\end{tcolorbox}
\section[\texorpdfstring{\ru{бедность} (bédnost)}{бедность (bédnost)}]{\ru{бедность} (bédnost) — poverty}

\label{lesson:RU-C367-bednost}



\begin{quote}
Before the new one: say the Russian for freedom, then the Russian for equality.
\end{quote}

\subsection*{You'll want to know: \ru{бедность}}
\textbf{\ru{бедность}} (\emph{bédnost}) — \textquotedblleft{}poverty\textquotedblright{}.

\begin{quote}
\ru{Многие} \ru{семьи} \ru{живут} \ru{в} \ru{бедности}.
\end{quote}

\begin{quote}
\emph{Mnógiye sémi zhivút v bédnosti.}
\end{quote}

\begin{quote}
Many families live in poverty.
\end{quote}

Feminine. From \emph{\ru{бедный}}, poor.

\begin{cousinweb}
\textbf{\ru{бедность}} is built from \emph{\ru{бедный}} \textquotedblleft{}poor\textquotedblright{} and the suffix \emph{-\ru{ость}} (EtymDB); \emph{\ru{бедный}} comes from \emph{\ru{беда}} \textquotedblleft{}trouble\textquotedblright{}, Proto-Slavic \emph{*běda} (EtymDB).
\end{cousinweb}

\begin{grammarlens}[title={What you've built}]
One more word for telling, explaining and describing at B1.
\end{grammarlens}

\subsection*{Your turn}
\begin{itemize}
\raggedright
  \item \emph{Say it:} \emph{bédnost}
  \item \emph{Say it:} \emph{bédnost}, once more
  \item \emph{Say it:} a sentence of your own with \emph{bédnost}
  \item \emph{Recall:} say the Russian for freedom, then the Russian for equality, then say \emph{bédnost} again
\end{itemize}

\begin{tcolorbox}[breakable,colback=teal!4,colframe=teal!35!black,title={Before you move on}]
What does \ru{бедность} mean? (\textquotedblleft{}Poverty\textquotedblright{}.) Say it once more.
\end{tcolorbox}
\section[\texorpdfstring{\ru{богатство} (bogátstvo)}{богатство (bogátstvo)}]{\ru{богатство} (bogátstvo) — wealth, riches}

\label{lesson:RU-C367-bogatstvo}



\begin{quote}
Before the new one: say the Russian for equality, then the Russian for poverty.
\end{quote}

\subsection*{You'll want to know: \ru{богатство}}
\textbf{\ru{богатство}} (\emph{bogátstvo}) — \textquotedblleft{}wealth, riches\textquotedblright{}.

\begin{quote}
\ru{Здоровье} — \ru{главное} \ru{богатство}.
\end{quote}

\begin{quote}
\emph{Zdoróvye — glávnoye bogátstvo.}
\end{quote}

\begin{quote}
Health is the greatest wealth.
\end{quote}

Neuter. \emph{\ru{Природные} \ru{богатства}} — natural resources.

\begin{cousinweb}
\textbf{\ru{богатство}} is built on \emph{\ru{богатый}} \textquotedblleft{}rich\textquotedblright{}, which is usually traced to \emph{\ru{бог}} in its old sense \textquotedblleft{}a share, good fortune\textquotedblright{}.
\end{cousinweb}

\begin{grammarlens}[title={What you've built}]
One more word for telling, explaining and describing at B1.
\end{grammarlens}

\subsection*{Your turn}
\begin{itemize}
\raggedright
  \item \emph{Say it:} \emph{bogátstvo}
  \item \emph{Say it:} \emph{bogátstvo}, once more
  \item \emph{Say it:} a sentence of your own with \emph{bogátstvo}
  \item \emph{Recall:} say the Russian for equality, then the Russian for poverty, then say \emph{bogátstvo} again
\end{itemize}

\begin{tcolorbox}[breakable,colback=teal!4,colframe=teal!35!black,title={Before you move on}]
What does \ru{богатство} mean? (\textquotedblleft{}Wealth, riches\textquotedblright{}.) Say it once more.
\end{tcolorbox}
\section[\texorpdfstring{\ru{безработица} (bezrabótitsa)}{безработица (bezrabótitsa)}]{\ru{безработица} (bezrabótitsa) — unemployment}

\label{lesson:RU-C367-bezrabotica}



\begin{quote}
Before the new one: say the Russian for poverty, then the Russian for wealth.
\end{quote}

\subsection*{You'll want to know: \ru{безработица}}
\textbf{\ru{безработица}} (\emph{bezrabótitsa}) — \textquotedblleft{}unemployment\textquotedblright{}.

\begin{quote}
\ru{В} \ru{этом} \ru{районе} \ru{высокая} \ru{безработица}.
\end{quote}

\begin{quote}
\emph{V étom rayóne vysókaya bezrabótitsa.}
\end{quote}

\begin{quote}
Unemployment is high in this area.
\end{quote}

Feminine. A person out of work is \emph{\ru{безработный}}.

\begin{cousinweb}
\textbf{\ru{безработица}} is built from \emph{\ru{без}} \textquotedblleft{}without\textquotedblright{} and \emph{\ru{работа}} \textquotedblleft{}work\textquotedblright{}, which goes back to Proto-Slavic \emph{*orbota} \textquotedblleft{}toil\textquotedblright{}, on \emph{\ru{раб}} \textquotedblleft{}a slave\textquotedblright{}.
\end{cousinweb}

\begin{grammarlens}[title={What you've built}]
One more word for telling, explaining and describing at B1.
\end{grammarlens}

\subsection*{Your turn}
\begin{itemize}
\raggedright
  \item \emph{Say it:} \emph{bezrabótitsa}
  \item \emph{Say it:} \emph{bezrabótitsa}, once more
  \item \emph{Say it:} a sentence of your own with \emph{bezrabótitsa}
  \item \emph{Recall:} say the Russian for poverty, then the Russian for wealth, then say \emph{bezrabótitsa} again
\end{itemize}

\begin{tcolorbox}[breakable,colback=teal!4,colframe=teal!35!black,title={Before you move on}]
What does \ru{безработица} mean? (\textquotedblleft{}Unemployment\textquotedblright{}.) Say it once more.
\end{tcolorbox}
